\ifx\allfiles\undefined
\documentclass[12pt, a4paper, oneside, UTF8]{ctexbook}
\def\path{../config}
\input{\path/package.tex}

% 定理环境设置
\input{\path/theorem1_zh.tex}

\input{\path/custom.tex}

% 修改参数改变封面样式，0 默认原始封面、内置其他1、2、3种封面样式
\def\myIndex{3}


\ifnum\myIndex>0
    \input{\path/cover_package_\myIndex} 
\fi

\def\myTitle{复变函数与积分变换笔记}
\def\myAuthor{M.K.}
\def\myDateCover{\today}
\def\myDateForeword{\today}
\def\myForeword{前言}
\def\myForewordText{
    
    这是一个基于\LaTeX{}模板的复变函数与积分变换笔记．
    内容参考了包括但不限于以下资料：
    \begin{itemize}
        \item 《复变函数与积分变换》（科学出版社）第三版
        \item 《复变函数习题精选精讲》（山东科学技术出版社）第一版
    \end{itemize}

    封面来源于\href{https://www.pixiv.net/users/2642047}{アシマ / Ashima}的作品\href{https://www.pixiv.net/artworks/147092236}{Assembly}
    如有侵权请联系我删除．
}
\def\mySubheading{副标题}


\begin{document}
%\input{../config/cover}
\else
\fi
\chapter{洛朗级数与孤立奇点}
\section{洛朗级数}
\begin{definition}[洛朗级数]
形如
\[
\sum_{n=-\infty}^{\infty} c_n (z-z_0)^n \defeq \sum_{n=0}^{\infty} c_n(z-z_0)^n + \sum_{n=1}^{\infty}\frac{c_{-n}}{(z-z_0)^n}
\]
的双边幂级数称为\textbf{Laurent 级数}，其中 $z_0\in\C$ 称为\textbf{中心}，$c_n\in\C$ 为系数。
\end{definition}

\begin{theorem}[Laurent 定理]
设 $D = \{z\in\C\mid r < \abs{z-z_0} < R\}$（$0\le r<R\le+\infty$）为环形区域，$f$ 在 $D$ 内解析。则 $f$ 在 $D$ 内可唯一展开为
\[
f(z) = \sum_{n=-\infty}^{\infty} c_n (z-z_0)^n,\quad z\in D,
\]
其中系数
\[
c_n = \frac{1}{2\pi i}\oint_\gamma \frac{f(z)}{(z-z_0)^{n+1}}\,\d z,\quad n\in\Z,
\]
$\gamma$ 为收敛圆环 $D$ 内任意绕 $z_0$ 的正向逐段光滑闭曲线。
\end{theorem}

\begin{example}
    求洛朗级数 $\displaystyle\sum_{n=1}^{\infty}\frac{2^n}{(z-3)^n}+\sum_{n=0}^{\infty}(-1)^n\left(1-\frac{z}{3}\right)^n$ 的收敛环。
\end{example}
\begin{solution}
    由 
    
    $\displaystyle\sum_{n=1}^{\infty}\frac{2^n}{(z-3)^n}$ 收敛当且仅当 $\abs{z-3}>2$
    
    $\displaystyle\sum_{n=0}^{\infty}(-1)^n\left(1-\frac{z}{3}\right)^n$ 收敛当且仅当 $\abs{1-\frac{z}{3}}<1$，即 $\abs{z-3}<3$

    因此原级数收敛环为 $2<\abs{z-3}<3$。
\end{solution}
\begin{example}
    求 $f(z)=z^2\e^{\frac{1}{z}}$ 在 $z_0=0$ 处的洛朗展开式。
\end{example}
\begin{solution}
    由 $\e^{\frac{1}{z}} = \sum_{n=0}^{\infty}\frac{1}{n!}\left(\frac{1}{z}\right)^n,\;\abs{z}>0$，得
    \[
    f(z) = z^2\e^{\frac{1}{z}} = z^2\sum_{n=0}^{\infty}\frac{1}{n!}\left(\frac{1}{z}\right)^n = \sum_{n=0}^{\infty}\frac{1}{n!}z^{2-n}.
    \]
    因此，$f(z)$ 在 $z_0=0$ 处的洛朗展开式为
    \[
    f(z) = \sum_{n=-2}^{\infty} \frac{1}{(n+2)!}z^n
    \]
\end{solution}
\begin{example}
    求 $f(z)=\dfrac{1}{(z^2+1)(z-2)}$ 在 $1<\abs{z}<2$ 处以 $z_0=0$ 为中心的洛朗展开式。
\end{example}
\begin{solution}
先分解为部分分式：
\[
\frac{1}{(z^2+1)(z-2)}
= -\frac{1}{5}\frac{z+2}{z^2+1} + \frac{1}{5}\frac{1}{z-2}.
\]

在 \(1<|z|<2\) 内分别展开：

\begin{itemize}
    \item 对于 \(\dfrac{1}{z-2}\)，因 \(|z|<2\)，有
        \[
        \frac{1}{z-2} = -\frac{1}{2}\cdot \frac{1}{1-\frac{z}{2}}
        = -\sum_{n=0}^{\infty} \frac{z^n}{2^{n+1}}.
        \]
    \item 对于 \(\dfrac{z+2}{z^2+1}\)，因 \(|z|>1\)，有
        \begin{align*}
            \frac{z+2}{z^2+1}
            &= \frac{1}{z}\left(1+\frac{2}{z}\right)\frac{1}{1+\frac{1}{z^2}}\\
            &= \frac{1}{z}\left(1+\frac{2}{z}\right)\sum_{k=0}^{\infty}(-1)^k z^{-2k}\\
            &= \sum_{k=0}^{\infty} (-1)^k z^{-2k-1} + 2\sum_{k=0}^{\infty} (-1)^k z^{-2k-2}
        \end{align*}
\end{itemize}
    合并得到
    \[
    \boxed{
    f(z)=\sum_{k=0}^{\infty}\frac{(-1)^{k+1}}{5}z^{-2k-1}
    +\sum_{k=0}^{\infty}\frac{2(-1)^{k+1}}{5}z^{-2k-2}
    -\sum_{n=0}^{\infty}\frac{z^n}{5\cdot 2^{n+1}},\quad 1<|z|<2.
    }
    \]
\end{solution}
\section{孤立奇点}

\subsection{奇点的分类}
\begin{definition}[孤立奇点]
设 $f$ 在 $z_0$ 的某去心邻域 $0<\abs{z-z_0}<\delta$ 内解析，但在 $z_0$ 处不解析，则称 $z_0$ 为 $f$ 的\textbf{孤立奇点}。否则称 $z_0$ 为\textbf{非孤立奇点}。
\end{definition}

\begin{definition}[奇点类型的分类]
设 $z_0$ 为 $f$ 的孤立奇点，$f$ 在 $0<\abs{z-z_0}<\delta$ 内的 Laurent 展开为
\[
f(z) = \sum_{n=-\infty}^{\infty} c_n (z-z_0)^n = \sum_{n=0}^{\infty} c_n(z-z_0)^n + \sum_{n=1}^{\infty}\frac{c_{-n}}{(z-z_0)^n}.
\]
称
\[
f_1(z) \defeq \sum_{n=0}^{\infty} c_n(z-z_0)^n
\]
为 $f$ 在 $z_0$ 处的\textbf{解析部分}（或\textbf{全纯部分}），
\[
f_2(z) \defeq \sum_{n=1}^{\infty}\frac{c_{-n}}{(z-z_0)^n}
\]
为 $f$ 在 $z_0$ 处的\textbf{主要部分}（或\textbf{奇异部分}）。按主要部分分类：
\begin{enumerate}
    \item 若 $f_2\equiv 0$（即所有 $c_{-n}=0$，$n\ge 1$），则称 $z_0$ 为 $f$ 的\textbf{可去奇点}；
    \item 若 $f_2$ 仅含有限项，设最高负幂项为 $\dfrac{c_{-m}}{(z-z_0)^m}$，$c_{-m}\neq 0$，$m\ge 1$，则称 $z_0$ 为 $f$ 的\textbf{$m$ 阶极点}；特别地 $m=1$ 时称\textbf{单极点}（或\textbf{简单极点}）；
    \item 若 $f_2$ 含无穷多项，则称 $z_0$ 为 $f$ 的\textbf{本性奇点}。
\end{enumerate}
\end{definition}

\subsection{奇点的判定}

\begin{theorem}[可去奇点的判别法]
设 $z_0$ 为 $f$ 的孤立奇点。下列命题等价：
\begin{enumerate}
    \item $z_0$ 为 $f$ 的可去奇点；
    \item $\displaystyle\lim_{z\to z_0} f(z)$ 存在（有限）；
    \item $f$ 在 $z_0$ 的某去心邻域内有界；
    \item Laurent 展开中所有负幂项系数均为零。
\end{enumerate}
\end{theorem}

\begin{theorem}[极点的判别法]
设 $z_0$ 为 $f$ 的孤立奇点。下列命题等价：
\begin{enumerate}
    \item $z_0$ 为 $f$ 的 $m$ 阶极点；
    \item $\displaystyle\lim_{z\to z_0} (z-z_0)^m f(z) = A\neq 0$（有限非零）；
    \item $\displaystyle\lim_{z\to z_0} f(z) = \infty$；
    \item $g(z)\defeq \dfrac{1}{f(z)}$ 以 $z_0$ 为 $m$ 阶零点（即可去奇点，且 $g^{(m)}(z_0)\neq 0$）。
\end{enumerate}
\end{theorem}

\begin{theorem}[本性奇点的判别法]
设 $z_0$ 为 $f$ 的孤立奇点。下列命题等价：
\begin{enumerate}
    \item $z_0$ 为 $f$ 的本性奇点；
    \item Laurent 展开中负幂项系数有无穷多个不为零；
    \item $\displaystyle\lim_{z\to z_0} f(z)$ 不存在且不为 $\infty$。
\end{enumerate}
\end{theorem}

\section{无穷远点}
\begin{definition}[无穷远点为孤立奇点]
    若函数 $f(z)$ 在 $z=\infty$ 的某去心邻域 $\{z\in\C\mid \abs{z}>R\}$ 内解析，则称 $f$ 以 $z=\infty$ 为孤立奇点。
\end{definition}
\begin{theorem}[无穷远点的奇点类型]
    设点 $\infty$ 为 $f$ 的孤立奇点，利用倒代换 $w = \frac{1}{z}$，将 $\infty$ 映为 $0$，于是 $g(w)=f(w^{-1})$ 在 $w=0$ 处有孤立奇点。则 $\infty$ 的奇点类型与 $0$ 的奇点类型相同。
\end{theorem}

\ifx\allfiles\undefined
\end{document}
\fi